%=============================================================================!
% 6.4  CONSTRAINED MOTION                                                     !
%=============================================================================!
\section{Constrained motion}
\label{sec:la-constrained}

The Euler--Lagrange equation holds in any generalised coordinates. Applied to the bead's single coordinate $x$, it
gives the equation of motion along the wire without the wire's push ever
appearing, and with it the times that \cref{sec:en-track} left to this
chapter.

\subsection*{The bead on the wire}

From \cref{sec:la-coordinates}, $\Lag = \frac12 m(1 + h'^2)\dot{x}^2 -
mgh(x)$, with $h$, $h'$ and $h''$ all functions of $x$. The two partial
derivatives are
\begin{equation*}
   \frac{\partial\Lag}{\partial\dot{x}} = m\,(1 + h'^2)\,\dot{x} ,
   \qquad
   \frac{\partial\Lag}{\partial x} = m\,h'h''\,\dot{x}^2 - mg\,h' ,
\end{equation*}
the second by the chain rule, with $\dot{x}$ held fixed, since
$\dd(h'^2)/\dd x = 2h'h''$. The rate
of change of the first along the motion needs the product rule, and the
chain rule for $h'^2$, which changes at the rate $2h'h''\dot{x}$:
\begin{equation*}
   \frac{\dd}{\dd t}\frac{\partial\Lag}{\partial\dot{x}}
   = m\,(1 + h'^2)\,\ddot{x} + 2m\,h'h''\,\dot{x}^2 .
\end{equation*}
Setting the two equal, by \cref{eq:la-euler},
\begin{align*}
   m\,(1 + h'^2)\,\ddot{x} + 2m\,h'h''\,\dot{x}^2
   &= m\,h'h''\,\dot{x}^2 - mg\,h'\\
   (1 + h'^2)\,\ddot{x} &= -h'h''\,\dot{x}^2 - g\,h'
   && \text{dividing by $m$, moving $2h'h''\dot{x}^2$ across}\\
   &= -h'\,\bigl(g + h''\dot{x}^2\bigr)
   && \text{taking out $-h'$,}
\end{align*}
and dividing by $1 + h'^2$,
\begin{equation}
   \ddot{x} = -\frac{h'\,\bigl(g + h''\dot{x}^2\bigr)}{1 + h'^2} .
   \label{eq:la-bead}
\end{equation}
The mass has cancelled. On a straight slope at the angle $\alpha$, $h =
x\tan\alpha$, $h' = \tan\alpha$ and $h'' = 0$. Dividing $\cos^2\alpha +
\sin^2\alpha = 1$ by $\cos^2\alpha$ gives $1 + \tan^2\alpha =
1/\cos^2\alpha$, so \cref{eq:la-bead} gives $\ddot{x} = -g\tan\alpha
\cos^2\alpha = -g\sin\alpha\cos\alpha$, with $\tan\alpha = \sin\alpha/
\cos\alpha$. The distance along the slope is the hypotenuse of a
right-angled triangle with horizontal side $x$, $x/\cos\alpha$, so its
acceleration is $\ddot{x}/\cos\alpha = -g\sin\alpha$: the part
$mg\sin\alpha$ of the weight down the ramp of \cref{sec:en-friction},
divided by $m$, with no friction. On a curved wire
the term in $h''\dot{x}^2$ adds the effect of the bending.

Why did the push of the wire never appear? It acts at right angles to the
wire, so it does no work in any motion along the wire; it has no
potential energy, and $\Lag = K - U$, made from the bead's speed and
height alone, does not contain it. Describing the bead by $x$ alone compares only paths
that keep it on the wire, and Newton's laws show that Hamilton's principle
holds among those. In $x$ and $y$ the bead obeys $m\ddot{\vb{r}} = -\grad
U + \vb{N}$, with $\vb{N}$ the push. Change the path by $\zeta
\boldsymbol\eta(t)$ along the wire, zero at both ends. The steps of
\cref{sec:la-euler}, once for $x$ and once for $y$, change the action by
$\zeta\int(-\grad U - m\ddot{\vb{r}})\cdot\boldsymbol\eta\,\dd t =
-\zeta\int\vb{N}\cdot\boldsymbol\eta\,\dd t$ plus terms in $\zeta^2$, and
this is zero because $\vb{N}$ is at right angles to the wire and
$\boldsymbol\eta$ along it. The true path therefore makes the action
stationary among the paths that stay on the wire, which are the paths of
$x$ alone, and the push never enters. \Cref{sec:la-multipliers} runs the
same argument backwards for the pendulum. This is the reason
for writing mechanics this way: constraints that do no work are obeyed by
choosing the coordinates, and their forces drop out.

\subsection*{The bead against the clock}

The function \texttt{solve\_wire} of \texttt{mdlab} solves
\cref{eq:la-bead}. Released from rest on the wire of \cref{sec:en-track}
at its left turning point for $h_E = \SI{1.3}{\metre}$, $x =
\SI{-2.206}{\metre}$ (\cref{sec:en-diagrams}), the bead
reaches the bottom of the deep valley after \SI{0.554}{\second}, the top
of the hill after \SI{1.097}{\second} and the bottom of the shallow valley
after \SI{1.533}{\second}, and turns at $x = \SI{2}{\metre}$ after
\SI{2.016}{\second}, coming back to its start after \SI{4.032}{\second}
(\cref{fig:la-bead}). Its speed along the wire, $\sqrt{1 + h'^2}\,
\lvert\dot{x}\rvert$, equals $\sqrt{2g(1.3 - h)}$ at every point, the
speed that \cref{sec:en-track} found from energy alone: 4.680, 2.387 and
\SI{3.686}{\metre\per\second} at the three places. The same times follow
from the energy, by adding up the time $\sqrt{1 + h'^2}\,\dd x/\sqrt{2g(1.3
- h)}$, length over speed, of each short stretch of wire, as
\cref{sec:os-anharmonic} did for the period; the equation of motion is the
route that still works with more than one coordinate.

\begin{figure}[htb]
\centering
\includegraphics{ch06_lagrange/bead}
\caption{The bead on the wire of \cref{sec:en-track}, released from rest
at $x = \SI{-2.206}{\metre}$, \SI{1.3}{\metre} high, its motion found from
\cref{eq:la-bead}. (a) Its position along $x$ against time, through one
swing to $x = \SI{2}{\metre}$ (dotted) and back; the dots mark the bottom
of the deep valley, the top of the hill and the bottom of the shallow
valley on the way out. (b) Its speed along the wire against $x$ on the way
out, lying on the speed from energy (dashed), which it hides.}
\label{fig:la-bead}
\end{figure}

\subsection*{The pendulum}

With $\Lag = \frac12 mL^2\dot\theta^2 - mgL(1 - \cos\theta)$,
$\partial\Lag/\partial\dot\theta = mL^2\dot\theta$ and $\partial\Lag/
\partial\theta = -mgL\sin\theta$, so \cref{eq:la-euler} gives
$mL^2\ddot\theta = -mgL\sin\theta$, or
\begin{equation}
   \ddot\theta = -\frac{g}{L}\sin\theta ,
   \label{eq:la-pendulum}
\end{equation}
in one line, with the rod's pull nowhere. For small swings $\sin\theta
\approx \theta$ (\cref{sec:mo-taylor}), and $\ddot\theta = -(g/L)\theta$
is the spring of \cref{sec:ne-spring} with $\omega_0 = \sqrt{g/L}$, as
\cref{sec:os-small} found from the energy.

\begin{keyidea}
Constraints that do no work are obeyed by choosing generalised
coordinates, and the Euler--Lagrange equations in those coordinates give
the motion without the constraining forces: for the bead, $\ddot{x} =
-h'(g + h''\dot{x}^2)/(1 + h'^2)$.
\end{keyidea}

\notebook{06}{put the bead on any wire and time its journey}

\begin{selfcheck}
How long does the bead released at \SI{1.3}{\metre} take to go from the
bottom of the deep valley to the bottom of the shallow one?
\end{selfcheck}
